\section{Worked RTL example and matched RF versus SFT evaluation}
\label{sec:supp-demo}
\subsection{What the example demonstrates}
The accompanying runnable package connects the paper's training objective to
an inspectable hardware example and then to the complete policy comparison.
The example uses \revisionclaim{DemoFirName}, the FIR design also discussed in
the main paper, with fresh outputs from the completed post-primary campaign.
For each policy, we select the most frequent functionally passing exact RTL
string among its \revisionclaim{DemoDrawsPerDesign} draws on this design,
breaking ties by the content hash. This selection is made after the experiment
to explain the implementation structure; timing is not used to choose either
string. The pair illustrates what a physical reward can favor, while the
all-design analysis below measures whether the resulting distribution is
better under the stated endpoint.

\begin{figure}[t]
\centering
\includegraphics[width=\columnwidth]{reviewer_demo.pdf}
\caption{Abbreviated assignments from the packaged FIR implementations.
The SFT output sums products of delayed inputs; the RF implementation registers
partial weighted sums. Full modules, including widths and reset logic, are
provided in the runnable package. These candidates illustrate a structural
difference and are not a causal ablation.}
\label{fig:reviewer-demo}
\end{figure}

In the SFT implementation, registers retain delayed inputs and a combinational
expression forms their weighted sum before the output register. In the RF
implementation, each stage stores a partial weighted sum, giving a transposed
FIR structure (Fig.~\ref{fig:reviewer-demo}). The selected strings occur in
\revisionclaim{DemoSftFirMultiplicity} SFT draws and
\revisionclaim{DemoRfFirMultiplicity} RF draws. Their measured bounded
setup-period frequencies are \revisionclaim{DemoSftFirFrequency} and
\revisionclaim{DemoRfFirFrequency}~MHz, respectively. Both pass the same
simulation gate, including repeated reset episodes and latency alignment.
This concrete difference is consistent with the intended reward, but does not
establish that the RF policy discovered a circuit absent from SFT support.
The measurements belong to the new campaign and do not replace the earlier
FIR measurements reported in the main paper.

\subsection{How to inspect and run the package}
After installing the dependency listed in \texttt{requirements.txt}, running
\texttt{python run\_demo.py} verifies the packaged hashes and recomputes the
comparison from all archived draws. The output includes the policy means,
functional pass counts, paired interval, per-design scores, and measured
selection curve. Adding \texttt{--simulate} runs the actual files
\texttt{demo/sft.sv} and \texttt{demo/rf.sv} with Icarus Verilog against the
shared reference. The gate uses \revisionclaim{DemoStreams} stimulus streams
of \revisionclaim{DemoStimuli} inputs, repeats each stream after reset, and
requires a common valid latency up to \revisionclaim{DemoMaxLatency} cycles
after a \revisionclaim{DemoWarmup}-sample warmup. Simulation traces and reset
checks are retained in the output. This is a functional simulation check;
it does not provide formal equivalence.

The package's \texttt{README.txt} provides installation commands, expected
values, file descriptions, and contact information. Candidate hashes link
the example to its exact draw and archived physical record; the selected
passing-boundary implementation reports are included for inspection. The
analysis runs without a GPU or Vivado because it replays archived physical
measurements. It does not retrain the policies or rerun implementation, and
the model weights and complete physical trial directories are outside the
package. These limits make the scope of the runnable demonstration explicit
while allowing the numerical comparison and the functional example to be
checked independently of the authors' server environment.

\subsection{Completed comparison on fresh draws}
% AUTO-GENERATED by paper/make_revision_evidence.py.
\begin{table}[t]
\centering
\caption{Completed post-primary RF versus SFT comparison. Utility is the all-draw bounded routed setup-period frequency, with failures scored as zero.}
\label{tab:demo-comparison}
\small
\begin{tabular}{@{}lrrr@{}}
\toprule
Policy & Passing draws & Pass (\%) & Utility (MHz) \\
\midrule
SFT & \revisionclaim{DemoSftPassed}/\revisionclaim{DemoDraws} & \revisionclaim{DemoSftRate} & \revisionclaim{DemoSftMean} \\
RF & \revisionclaim{DemoRfPassed}/\revisionclaim{DemoDraws} & \revisionclaim{DemoRfRate} & \revisionclaim{DemoRfMean} \\
\bottomrule
\end{tabular}

\end{table}

The completed campaign evaluates \revisionclaim{DemoDesigns} previously used
designs with \revisionclaim{DemoDrawsPerDesign} new draws per policy and design.
SFT has two sampling replicates from the same checkpoint; RF uses the existing
training-seed-three and training-seed-four checkpoints, each generating
\revisionclaim{DemoReplicateDraws} draws per design. All outputs use the shared
functional gate and bounded routed setup-period search with Vivado
\revisionclaim{DemoVersion}. Every functional or physical failure contributes
zero, and each design receives equal weight. This post-primary comparison
comes from a completed campaign that also evaluated a DPO arm. The present
package isolates the RF versus SFT contrast requested for the supplement;
it is not a matched reproduction of an external method or an independent
test set, and its measurements are not pooled with the primary study.

RF reaches \revisionclaim{DemoRfMean}~MHz in failure-aware utility, compared
with \revisionclaim{DemoSftMean}~MHz for SFT
(Table~\ref{tab:demo-comparison}). The paired difference is
\revisionclaim{DemoGain}~MHz, with a \revisionclaim{DemoConfidence}\%
family-by-regime design-bootstrap interval of
[\revisionclaim{DemoLow}, \revisionclaim{DemoHigh}]~MHz from
\revisionclaim{DemoResamples} resamples. The RF replicate means are
\revisionclaim{DemoRfFirst} and \revisionclaim{DemoRfSecond}~MHz, compared
with \revisionclaim{DemoSftFirst} and \revisionclaim{DemoSftSecond}~MHz for
the SFT sampling replicates. RF's functional pass rate is lower:
\revisionclaim{DemoRfRate}\% versus \revisionclaim{DemoSftRate}\%.
The physical utility gain therefore coexists with a correctness tradeoff.
The interval is conditional on these checkpoints and draws and does not
capture uncertainty from retraining SFT or the entire pipeline. The endpoint
measures bounded routed setup-period performance, not full timing signoff.

\subsection{Selection at equal draw budgets}
% AUTO-GENERATED by paper/make_revision_evidence.py.
\begin{table}[t]
\centering
\caption{Measured best-of-$N$ utility in MHz, using fixed consecutive blocks within each generation replicate. Selection uses measured physical outcomes.}
\label{tab:demo-selection}
\small
\begin{tabular}{@{}rrr@{}}
\toprule
$N$ & SFT & RF \\
\midrule
\revisionclaim{DemoBudgetOne} & \revisionclaim{DemoSftBestOne} & \revisionclaim{DemoRfBestOne} \\
\revisionclaim{DemoBudgetFour} & \revisionclaim{DemoSftBestFour} & \revisionclaim{DemoRfBestFour} \\
\revisionclaim{DemoBudgetEight} & \revisionclaim{DemoSftBestEight} & \revisionclaim{DemoRfBestEight} \\
\revisionclaim{DemoBudgetTwentyFour} & \revisionclaim{DemoSftBestTwentyFour} & \revisionclaim{DemoRfBestTwentyFour} \\
\bottomrule
\end{tabular}

\end{table}

To separate policy quality from the opportunity to select among several
outputs, Table~\ref{tab:demo-selection} reports the best measured result in
fixed consecutive draw blocks within each replicate. Failures remain zero,
and the same draw budgets are used for both policies. RF has higher utility
at each reported budget, while the gap narrows as more candidates are
available. At the largest budget, SFT reaches
\revisionclaim{DemoSftBestTwentyFour}~MHz and RF reaches
\revisionclaim{DemoRfBestTwentyFour}~MHz. This comparison grants the selector
access to measured physical outcomes; it does not model an inexpensive
reranker or include training cost. Together, the example and the complete
draw-level comparison support a shift toward physically faster outputs under
the stated evaluation contract, while leaving the overall compute tradeoff
and transfer to broader RTL workloads open.
